\begin{proofbox}
\textbf{统一记号：}
\[
F_1(-c,0),\quad F_2(c,0),\quad
r_1=PF_1,\quad r_2=PF_2,\quad
\angle F_1PF_2=\theta.
\]
由椭圆定义与焦距公式，
\[
r_1+r_2=2a,\qquad F_1F_2=2c,\qquad b^2=a^2-c^2.
\]

\medskip
\textbf{证明一：几何投影法（最直观，建议先掌握）。}

\begin{enumerate}
    \item 作 $\angle F_1PF_2$ 的内部角平分线 $\ell$，取一个与 $\ell$ 垂直的方向。
    \item 因为 $PF_1,PF_2$ 分居角平分线两侧，所以它们在该垂直方向上的投影方向相反，投影长度分别为
    $r_1\sin(\theta/2)$、$r_2\sin(\theta/2)$。
    \item 因此，焦点连线 $F_1F_2$ 在该方向的投影长度为
    \[
    r_1\sin\frac\theta2+r_2\sin\frac\theta2
    =(r_1+r_2)\sin\frac\theta2
    =2a\sin\frac\theta2.
    \]
    \item 任何线段在任意方向上的投影长度都不大于原线段长度，故
    \[
    2a\sin\frac\theta2\le F_1F_2=2c
    \quad\Longrightarrow\quad
    \boxed{\sin\frac\theta2\le\frac ca=e.}
    \]
    \item 等号成立当且仅当 $F_1F_2\perp\ell$。
    此时两焦点在角平分线方向上的投影相同，即
    $r_1\cos(\theta/2)=r_2\cos(\theta/2)$。
    由于 $0\le\theta<\pi$，故 $r_1=r_2$，从而 $x_P=0$、$P=(0,\pm b)$。
\end{enumerate}

\medskip
\textbf{证明二：余弦定理 $+$ 基本不等式（保留原证明并补足细节）。}

在焦点三角形中，由余弦定理
\begin{align*}
\cos\theta
&=\frac{r_1^2+r_2^2-(2c)^2}{2r_1r_2}\\
&=\frac{(r_1+r_2)^2-2r_1r_2-4c^2}{2r_1r_2}\\
&=\frac{4(a^2-c^2)-2r_1r_2}{2r_1r_2}
=\frac{2b^2}{r_1r_2}-1.
\end{align*}
再利用半角公式，
\[
\sin^2\frac\theta2
=\frac{1-\cos\theta}{2}
=1-\frac{b^2}{r_1r_2}.
\]
由 $r_1+r_2=2a$ 及基本不等式，
\[
r_1r_2\le\left(\frac{r_1+r_2}{2}\right)^2=a^2,
\]
所以
\[
\sin^2\frac\theta2
=1-\frac{b^2}{r_1r_2}
\le1-\frac{b^2}{a^2}
=\frac{c^2}{a^2}=e^2.
\]
因 $\sin(\theta/2)\ge0$、$e>0$，开方得
\[
\boxed{\sin\frac\theta2\le e.}
\]
等号成立当且仅当 $r_1=r_2=a$，也就是 $P=(0,\pm b)$。

\medskip
\textbf{证明三：余弦定理的半角分解（直接看出“差多少”）。}

从余弦定理出发，但不急着使用基本不等式：
\begin{align*}
(2c)^2
&=r_1^2+r_2^2-2r_1r_2\cos\theta\\
&=(r_1-r_2)^2+4r_1r_2\sin^2\frac\theta2\\
&=(r_1-r_2)^2\cos^2\frac\theta2
 +(r_1+r_2)^2\sin^2\frac\theta2\\
&=(r_1-r_2)^2\cos^2\frac\theta2
 +4a^2\sin^2\frac\theta2.
\end{align*}
于是得到\textbf{精确恒等式}
\[
\boxed{
 e^2-\sin^2\frac\theta2
 =\frac{(r_1-r_2)^2}{4a^2}\cos^2\frac\theta2\ge0.
}
\]
右侧非负，立即得证。它还说明：
\emph{在两焦半径不相等时，非零的“长度差”占用了焦距的一部分，角度就达不到最大值。}
等号只能在 $r_1=r_2$ 时成立。

\medskip
\textbf{证明四：焦半径公式 $+$ 坐标单调性（适合解析几何计算）。}

令 $P=(x,y)$。椭圆的焦半径公式为
\[
r_1=a+ex,\qquad r_2=a-ex.
\]
因此
\[
r_1r_2=a^2-e^2x^2.
\]
代入证明二中的半角表达式，得到
\[
\boxed{
\sin^2\frac\theta2
=1-\frac{b^2}{a^2-e^2x^2}.
}
\]
椭圆上 $|x|\le a$，分母满足
$ a^2-e^2x^2\ge a^2-c^2=b^2>0$。
随着 $|x|$ 增大，分母减小，分数 $b^2/(a^2-e^2x^2)$ 增大，
故 $\sin^2(\theta/2)$ 减小。
于是当且仅当 $x=0$ 时取最大值：
\[
\sin^2\frac{\theta_{\max}}2
=1-\frac{b^2}{a^2}=e^2.
\]
从而 $\theta_{\max}=2\arcsin e$，对应 $P=(0,\pm b)$。
反过来，$x=\pm a$ 时上式等于 $0$，故长轴端点的张角为 $0$。

\medskip
\textbf{结论的“当且仅当”为什么成立？}

已经证明每一点的 $\theta\le\theta_{\max}=2\arcsin e$。
沿椭圆的第一象限弧，从 $(a,0)$ 连续移动到 $(0,b)$，
张角连续地由 $0$ 变为 $\theta_{\max}$。
由连续性，每个 $0\le\theta_0\le\theta_{\max}$ 都能取到。
因此对 $0\le\theta_0<\pi$，
\[
\exists P:\ \angle F_1PF_2\ge\theta_0
\quad\Longleftrightarrow\quad
\theta_{\max}\ge\theta_0
\quad\Longleftrightarrow\quad
 e\ge\sin\frac{\theta_0}{2}.
\]
\end{proofbox}
